% DRAFT (round 14) -- new subsection after \subsection{Which Contexts Share a Concept}; wording of the purity
% claim depends on the verdict of rule P (strict / majority / conditional).

\subsection{Which Rows Form a Concept}
\label{sec:th-seg}

\Cref{prop:recall} compares relearning with the stored expert \emph{of concept $k$}. A stored context is built from
segments of the stream, and a segment that straddles a concept boundary holds rows of two concepts. The next two
results say how much mixing the recall argument tolerates and how a change detector bounds the mixing.

\begin{proposition}[Majority purity]\label{prop:purity}
Let \cref{ass:ppd} extend to unequal shares: a context whose rows come from $P$ in proportion $\pi$ and from $Q$ in
proportion $1-\pi$ predicts the conditional of $\pi P+(1-\pi)Q$. Let $P_X=Q_X$ and let both concepts be
deterministic, $y=g_P(\vx)$ and $y=g_Q(\vx)$. Under the 0-1 loss, the recall regret of this context on $P$ is
\[
  \Div_T(P\,|\,\pi P+(1-\pi)Q)=
  \begin{cases}
    0, & \pi>\tfrac12,\\
    \rho, & \pi<\tfrac12,
  \end{cases}
\]
where $\rho=\Pr_\vx[g_P(\vx)\ne g_Q(\vx)]$; at $\pi=\tfrac12$ it depends on the tie rule, as in
\cref{prop:blind}(iii).
\end{proposition}

\begin{remark}
In the population limit a stored context recalls the concept that holds the majority of its rows, so
identifying the concept needs majority purity, not purity. \Cref{prop:recall} then applies to the majority concept
with a finite-context correction: the minority rows are wrong labels for it on the disagreement region and correct
ones elsewhere, so the context of $n$ rows with share $\pi>\tfrac12$ is worth between $\pi n$ and $n$ rows of the
majority concept. Dropping the rows near a boundary raises $\pi$ but lowers $n$; whether the worth rises or falls
depends on the disagreement and the learning curve (\cref{sec:exp-seg}).
\end{remark}

\textbf{Change detection.} Let $e_t\in\{0,\dots,B\}$ be the number of errors that the $300$-row window makes on
batch $t$, scored when its labels arrive. For a candidate change at batch $b$, compare the error rates on the up to
$H$ batches after it, $b,\dots,t-1$ ($n_2$ rows), with those on the up to $R$ batches before it ($n_1$ rows), and
declare a change at the $b$ with the largest excess
\begin{equation}\label{eq:cp}
  \frac{1}{n_2}\sum_{s=b}^{t-1}e_s-\frac{1}{n_1}\sum_{s=b-R}^{b-1}e_s-\varepsilon(n_1,n_2),\qquad
  \varepsilon(n_1,n_2)=\sqrt{\tfrac12\bigl(\tfrac1{n_1}+\tfrac1{n_2}\bigr)\ln\tfrac1{\delta_{\mathrm{cp}}}},
\end{equation}
if that excess is positive. The compared batches before $b$ start after the batch that follows the previous
declared change, and at least two are required; $n_1$ and $n_2$ count $B$ rows per compared batch.

\begin{assumption}[Window error]\label{ass:err}
Given its context, the window errs on the rows of a batch independently, each with the conditional error rate of
that batch. (a) Between two changes this rate does not increase from batch to batch. (b) After a change at batch $b$,
the mean rate on the batches $b,\dots,t-1$ exceeds that on the compared batches before $b$ by at least
$\Delta$.
\end{assumption}

\begin{proposition}[Detection]\label{prop:cp}
Under \cref{ass:err}:
\begin{enumerate}
  \item[(i)] \emph{False alarms.} If no change occurs in the batches compared at time $t$, the probability that
  \eqref{eq:cp} declares one at $t$ is at most $H\delta_{\mathrm{cp}}$.
  \item[(ii)] \emph{Detection.} If a change occurs at batch $b$ and $\Delta\ge2\varepsilon(n_1,n_2)$ at time $t$,
  then at $t$ the excess at $b$ is positive, and a change is declared, with probability at least
  $1-\delta_{\mathrm{cp}}$.
\end{enumerate}
\end{proposition}

\begin{corollary}[Segment purity]\label{cor:cp}
If every declared change lies within $D$ rows of a true boundary and no boundary is missed, then a closed segment
of $L_S$ rows holds at most $2D$ rows of other concepts, and it is majority-pure when $4D<L_S$.
\end{corollary}

\begin{remark}
With $B=100$, $H=3$, $R=5$ and $\delta_{\mathrm{cp}}=10^{-3}$, the per-batch false-alarm probability is at
most $0.003$. By \cref{prop:purity}, the window, whose $300$ rows hold the previous concept in the majority for two
batches after a change, keeps predicting that concept on those batches; for two deterministic concepts that disagree
on mass $\rho$, the gap is about $\rho$ on batches $b$ and $b+1$. With $n_1=500$ and $n_2=200$,
$2\varepsilon=0.31$, so changes between concepts that disagree on more than about a third of the inputs are detected
within two batches with probability at least $0.999$; changes between closer concepts may be missed, and a
segment then holds both. When the declared batch is the one that holds the boundary or the next, $D<B$, and a segment of at least
$400$ rows is majority-pure; how often this holds is measured in \cref{sec:exp-seg}. Segmentation by detected changes therefore implements the premise
of \cref{prop:recall} up to the stated mixing, wherever the boundaries fall.
\end{remark}

% Proofs (appendix)
% prop:purity: the mixture conditional is pi*1[y=g_P] + (1-pi)*1[y=g_Q]; where g_P=g_Q the argmax is correct; where
%   they differ the argmax is g_P iff pi>1/2. 0-1 risk on P = (pi<1/2)*rho; risk of h_P on P is 0.
% prop:cp (i): the per-row error indicators minus their conditional means form a martingale difference sequence
%   (a batch's context is fixed before its rows are seen);
%   the statistic (post mean - pre mean) has increments bounded by 1/n2 and 1/n1, sum of squares 1/n1+1/n2;
%   Azuma-Hoeffding gives P(stat - E >= eps) <= exp(-2 eps^2/(1/n1+1/n2)) = delta_cp; E <= 0 without a change;
%   union bound over the <= H candidates b.  (ii): E >= Delta; P(stat <= eps) <= P(stat - E <= -(Delta-eps))
%   <= exp(-2 (Delta-eps)^2/(1/n1+1/n2)) <= delta_cp when Delta-eps >= eps.
% cor:cp: each end of the segment is within D rows of a boundary, so at most D foreign rows at each end; majority
%   if 2D < L_S/2.
